\section{Worked synthesis: identities, bytes and evidence}
\subsection{A vocabulary permutation that passes shape checks}
\textbf{Problem.} Two tokenizers have the same vocabulary size. One exchanges
the IDs for two ordinary tokens. Why can model shape validation succeed while
the answer changes?

\textbf{Solution.} A shape check proves only that an ID is in range. The same
integer now selects a different embedding row and labels a different output
logit. If only the tokenizer changes, the model's learned association between
row and token is broken. A coordinated permutation of embedding rows,
output rows, biases and all special-token semantics could preserve behavior,
but changing one file does not perform that transformation.

\subsection{Complete the byte trace}
\textbf{Problem.} A token emits \texttt{E4 B8}; the next emits \texttt{96}.
List pending bytes and text after each token. What changes if EOS arrives
instead of the second token?

\textbf{Solution.} After the first token, pending bytes are
\texttt{E4 B8} and text is empty. After the second, the pending buffer is empty
and U+4E16 has been emitted once. EOS after the first token leaves an incomplete
UTF-8 scalar; under this engine's strict completion policy it is a decoding
failure, not successful text and not an implicit replacement character.

\subsection{Build a discriminating tokenizer check}
\textbf{Problem.} Is encode/decode round-trip equality sufficient to validate
compatibility with a model?

\textbf{Solution.} No. A self-consistent but permuted vocabulary round-trips its
own IDs. Add known byte-to-ID fixtures from the intended model revision,
special-token checks and template-rendering examples. Round trips test internal
consistency; known fixtures test the external semantic contract.
